\documentclass[11pt,a4paper]{article}
\usepackage[margin=1in]{geometry}
\usepackage{amsmath}
\usepackage{enumitem}
\usepackage{siunitx}
\setlength{\parskip}{6pt}
\setlength{\parindent}{0pt}

\title{Cálculo de trayectorias y métricas en levantamientos verticales}
\author{MOOVIA Sensor Lab}
\date{Marzo 2026}

\begin{document}
\maketitle

\section*{Objetivo}
Describir paso a paso cómo el pipeline convierte las muestras crudas del IMU en aceleración lineal, velocidad, desplazamiento vertical y métricas (aceleración pico, velocidad media propulsiva y altura máxima) usadas en el panel \emph{Processed}.

\section{Entrada y reconstrucción temporal}
\begin{itemize}[leftmargin=12pt]
  \item Cada paquete BLE trae 15 muestras (\(ax, ay, az, gx, gy, gz, hwTs\)) a \SI{1}{kHz}.
  \item Para cada paquete, se reconstruye el tiempo absoluto:
  \[
    \Delta t = \frac{\Delta \text{ticks} \times 32/30\ \mu s}{1000}, \qquad
    t_i = t_{i+1} - \Delta t
  \]
  usando wrap-around de 16 bits. El último timestamp del paquete se ancla al reloj del sistema si hay huecos \(>\) \SI{1}{s}.
\end{itemize}

\section{Calibración inicial}
\begin{itemize}[leftmargin=12pt]
  \item Durante \(\approx\) \SI{2}{s} en reposo se promedian las aceleraciones crudas.
  \item Se construye un cuaternión \(q_0\) que alinea ese vector con el eje \(+Z\) del mundo (gravedad abajo). Se reinician cinemática y filtros.
\end{itemize}

\section{Transformación a marco mundo y compensación de gravedad}
\begin{enumerate}[leftmargin=14pt]
  \item Convertir aceleración a \(\text{m/s}^2\): \(a_b = (ax, ay, az)\cdot 9.81\).
  \item Rotar al marco mundo con el cuaternión actual \(q\) (sensor\(\rightarrow\)mundo):
  \[
    a_w = q \otimes a_b \otimes q^{-1}
  \]
  \item Estimar gravedad con filtro exponencial (coef. \(\alpha=0.98\)):
  \[
    g_{est} = \alpha\, g_{est} + (1-\alpha)\, a_w
  \]
  \item Aceleración lineal:
  \[
    a_{lin} = a_w - g_{est}
  \]
\end{enumerate}

\section{Integración con filtros Kalman 1D}
Para cada eje \(x,y,z\) se usa un estado \(x = [p,\ v,\ b]^T\) (posición, velocidad, sesgo de acelerómetro):
\begin{align*}
  \text{Predicción: } & 
  v_{k+1} = v_k + (a_{lin}-b_k)\,\Delta t,\\
  p_{k+1} = p_k + v_k\,\Delta t + \tfrac{1}{2}(a_{lin}-b_k)\Delta t^2.\\
  \text{Actualización (si está parado): } &
  v = 0,\quad b = a_{lin} \quad (\text{ZUPT + bias}).
\end{align*}
La detección de reposo usa ventanas deslizantes: \(|a|\) cerca de \(g\), \(\text{std}(a) < 0.15\ \text{m/s}^2\), \(\text{std}(\omega) < 0.05\ \text{rad/s}\) con histéresis (entrada 300 ms, salida 100 ms).

\section{Trayectoria y orientación para la UI}
\begin{itemize}[leftmargin=12pt]
  \item Se acumulan puntos \((t, p, v, q)\); para la vista frontal se aplica la rotación fija de montura y el ajuste \(-90^\circ\) en \(X\) (convertir mundo \(Z\uparrow\) a UI \(Y\uparrow\)).
  \item Al detener el stream se reprocesan todos los datos crudos para reconstruir la trayectoria relativa al primer punto (baseline).
\end{itemize}

\section{Métricas mostradas}
\begin{itemize}[leftmargin=12pt]
  \item \textbf{Aceleración Pico}: \(\max_t \|a_{lin}(t)\|\) en \(\text{m/s}^2\).
  \item \textbf{V. Media Propulsiva}: media de \(v_z(t)\) durante la fase concéntrica (\(v_z > 0.05\ \text{m/s}\)).
  \item \textbf{Altura Máx}: \(\max_t (p_z(t) - p_z(0))\) tras el reprocesado offline.
\end{itemize}

\section{Cómo reproducir el análisis}
\begin{enumerate}[leftmargin=14pt]
  \item Exporta la sesión desde la app móvil (JSON).
  \item Abre \texttt{apps/web-test} con Vite y pulsa \texttt{Load JSON}.
  \item La web usa el mismo pipeline descrito arriba para graficar aceleración, velocidad y desplazamiento vertical.
\end{enumerate}

\section*{Notas para levantamientos verticales}
\begin{itemize}[leftmargin=12pt]
  \item Asegura 1--2 s en reposo al inicio para que la calibración capture la gravedad correctamente.
  \item Evita huecos de transmisión: si \(\Delta t > 50\ \text{ms}\) se congela la integración hasta que vuelvan muestras estables.
  \item Revisa el gráfico de velocidad: un perfil limpio ascenso--descenso debe ser simétrico; picos espurios suelen venir de golpes o drift si no hubo reposo suficiente.
\end{itemize}

\end{document}
